\ncdlesson{1}{Et postkort fra Rom --- og en e-mail hjem}{来自罗马的明信片——以及一封寄回家的邮件}
{DU3 Modul 2（A2）｜对标 NCE 2 Lesson 3 \& Modul 2 \textit{Skrivning Opgave 1}}
{现在完成时（\textit{Førnutid}）助动词 \textbf{\textit{har} vs \textit{er}} · \textbf{现在完成时 vs 一般过去时（\textit{Datid}）切换法则} · Modul 2 邮件写作四步模板}
{NCD 2 每课采用“\textbf{经典故事短文（Tekst A）+ Modul 2 真题实战（Tekst B）}”双轨结构。冲刺 Modul 2 时，请重点精读 Tekst B 的回信结构与第四部分的时态切换法则——这两项是写作拿高分的直接分水岭。}

% ---------------------------------------------------------------------
\begin{foerduleaeser}
带着下面两个问题阅读 Tekst A 和 Tekst B：
\begin{enumerate}
  \item \textit{Hvor længe har Lukas og Maja været i Rom, og hvor mange postkort har Lukas skrevet indtil nu?}（卢卡斯和玛雅在罗马待了多久？卢卡斯目前为止写了几张明信片？）
  \item \textit{Hvad lavede de i går ifølge Majas e-mail, og hvordan har vejret været?}（根据玛雅的邮件，他们昨天做了什么？天气一直怎么样？）
\end{enumerate}
\end{foerduleaeser}

% ---------------------------------------------------------------------
\ncdsubsec{二、Tekst A \& Tekst B（双轨课文、发音提示与中文对照）}

\subsection*{Tekst A：\textit{Et postkort fra Rom}（叙事短文 · 约 155 词）}

\begin{tipbox}{语篇说明}
改编自《新概念英语 2》第 3 课 \textit{Please send me a card}。请特别留意文中\textbf{\color{NCDNavy}深蓝色加粗的现在完成时}（\textit{har/er + 分词}）与一般过去时的交替节奏。
\end{tipbox}

Lukas og hans kæreste, Maja, \da{er rejst} til Rom på ferie. De \da{har} allerede \da{været} i Italien i fem dage, og de \da{har oplevet} utrolig meget.

Før de tog afsted, sagde Lukas altid, at postkort ødelægger en god ferie, fordi man bruger alt for meget tid på at skrive dem. Alligevel gik han ind i en lille butik den første dag og købte tyve smukke postkort.

Hver formiddag \da{har} de \da{besøgt} berømte museer og gamle kirker. Om eftermiddagen \da{har} Lukas \da{siddet} på en hyggelig café i solen og \da{drukket} espresso. Han \da{har lagt} postkortene frem på bordet og \da{har prøvet} at skrive til sin familie og sine kolleger i Danmark. Men indtil nu \da{har} han ikke \da{skrevet} et eneste ord!

Maja \da{har} derimod allerede \da{sendt} mange billeder til sine venner på telefonen, og i går skrev hun en lang e-mail tilbage til sin ven Daniel. I morgen flyver de hjem til København, og Lukas tager nok de tyve tomme postkort med hjem i kufferten.

\begin{tcolorbox}[colback=NCDBgGray, colframe=NCDGray!40, boxrule=0.5pt, arc=2pt, breakable]
\small
\textbf{\color{NCDTeal}[Tekst A 中文精译]}\\
卢卡斯和他的女友玛雅去罗马度假了。他们在意大利已经待了五天，经历了非常多精彩的事情。

出发之前，卢卡斯总是说，明信片会毁掉一个好假期，因为人会花太多时间去写它们。尽管如此，第一天他还是走进一家小店，买了二十张漂亮的明信片。

每天上午，他们都去参观了著名的博物馆和古老的教堂。每天下午，卢卡斯都坐在阳光下一家惬意的咖啡馆里喝意式浓缩咖啡。他把明信片摆在桌面上，试着给丹麦的家人和同事写点什么。可是直到现在，他一个字也还没写出来！

相反，玛雅已经在手机上给朋友们发了好多照片，而且昨天她还给好朋友丹尼尔回了一封长长的电子邮件。明天他们就要飞回哥本哈根了，卢卡斯大概会把那二十张空白的明信片装在行李箱里带回家。
\end{tcolorbox}

\subsection*{Tekst B：\textit{En e-mail fra Daniel --- og Majas svar}（对标 DU3 Modul 2 Skrivning Opgave 1）}

\textbf{来信题目（Situation）：}\textit{Du har fået en e-mail fra din ven, Daniel. Læs e-mailen, og skriv et svar til Daniel. Du skal svare på alle Daniels spørgsmål. Du skal skrive minimum 70 ord.}

\begin{emailbox}{来信｜Fra: Daniel}
\textbf{Hej Maja}\\
Jeg har hørt, at du er ude at rejse. Det lyder spændende! Hvor er du henne? Hvem rejser du sammen med? Hvad laver I? Og hvordan er vejret?\\
\textbf{Mange hilsner}\\
\textbf{Daniel}
\end{emailbox}

\begin{emailbox}{高分回信范本（96 词）｜Til: Daniel \quad Emne: Hilsen fra Rom}
\textbf{Hej Daniel}

Tak for din søde mail! Ja, det er rigtigt, at jeg er ude at rejse. Jeg er i Rom sammen med min kæreste, Lukas. Vi ankom i mandags, så vi \da{har} allerede \da{været} her i fem dage.

Vi \da{har set} mange spændende museer og gamle kirker. I går gik vi en lang tur gennem byen, og om aftenen spiste vi god pizza på en lille restaurant. Vejret \da{har været} fantastisk med masser af sol, selvom det blæste lidt den første dag.

Hvordan går det derhjemme i København? Jeg glæder mig til at se dig, når vi kommer hjem på søndag!

\textbf{Kærlig hilsen}\\
\textbf{Maja}
\end{emailbox}

\begin{tipbox}{本课发音与吞音要点（Udtale \& Reduktioner）}
\begin{itemize}
  \item \textbf{过去分词词尾 \textit{-et}}：绝对不要读成硬邦邦的 \ipa{ɛt}！标准发音是轻弱的\textbf{央元音 + 软 d}：\ipa{əð}。例如：\textit{oplevet} \ipa{ˈʌbˌleːˀvəð}、\textit{været} \ipa{ˈvæɐ̯ð} / \ipa{ˈvɛːɐð}、\textit{skrevet} \ipa{ˈsgʁɛːvəð}、\textit{drukket} \ipa{ˈdʁɔgəð}。
  \item \textbf{助动词 \textit{har} 与 \textit{er} 的弱读}：在句子里作助动词时几乎完全不重读，发音像轻短的 \ipa{hɑ} 和 \ipa{æɐ̯}，重音落在后面的\textbf{主要动词分词}上：\textit{De har \textbf{været}} \ipa{di hɑ ˈvæɐ̯ð}，\textit{De er \textbf{rejst}} \ipa{di æɐ̯ ˈʁɑɪ̯ˀsd}。
  \item \textbf{连读吞音}：\textit{i fem dage} 中的 \textit{dage} 读 \ipa{ˈdæːə}（\textbf{g 不发音}）；\textit{altid} \ipa{ˈalˀtiðˀ}；\textit{i går} \ipa{i ˈɡɒːˀ}。
\end{itemize}
\end{tipbox}

\begin{tcolorbox}[colback=NCDBgGray, colframe=NCDGray!40, boxrule=0.5pt, arc=2pt, breakable]
\small
\textbf{\color{NCDTeal}[Tekst B 中文对照]}\\
\textbf{丹尼尔}：嗨，玛雅！我听说你出门旅行了。听起来真棒！你在哪儿呢？你和谁一起旅行？你们在做什么？天气怎么样？祝好，丹尼尔。

\textbf{玛雅回信}：嗨，丹尼尔！谢谢你亲切的邮件！是的，我确实在外面旅行。我和我男朋友卢卡斯在罗马。我们周一到的，所以已经在这儿待了五天了。我们看了很多有意思的博物馆和老教堂。昨天我们在城里散了很长一段步，晚上在一家小餐馆吃了美味的披萨。天气一直棒极了，阳光充足，虽然第一天有点刮风。哥本哈根家里那边怎么样？我真期待周日回家后见到你！亲切的问候，玛雅。
\end{tcolorbox}

% ---------------------------------------------------------------------
\ncdsubsec{三、Gloser（核心生词表）}

\begin{tipbox}{重点标记}
下面用 \dabold{[er]} 标出的动词，构成现在完成时必须用 \textit{er}，绝不能用 \textit{har}！
\end{tipbox}

{\small
\renewcommand{\arraystretch}{1.25}
\begin{tabularx}{\linewidth}{@{} >{\bfseries\color{NCDGray}}l >{\itshape\bfseries\color{NCDNavy}}l >{\raggedright\arraybackslash}p{6.5cm} >{\raggedright\arraybackslash}X @{}}
\toprule
\textbf{词性} & \textbf{丹麦语原形} & \textbf{完整屈折变化（含完成时助动词）} & \textbf{中文释义} \\
\midrule
动词 \dabold{[er]} & at rejse & \textit{(rejser, rejste, \dabold{er rejst})} & 旅行；出发前往 \\
动词 \dabold{[er]} & at ankomme & \textit{(ankommer, ankom, \dabold{er ankommet})} & 到达，抵达 \\
动词 \dabold{[er]} & at tage afsted & \textit{(tager, tog, \dabold{er taget afsted})} & 出发，动身 \\
动词 \dabold{[er]} & at flyve & \textit{(flyver, fløj, \dabold{er fløjet})} & 飞，乘飞机 \\
动词 \dabold{[er]} & at gå & \textit{(går, gik, \dabold{er gået} / har gået)} & 走；去（见语法辨析） \\
\midrule
动词 [har] & at opleve & \textit{(oplever, oplevede, har oplevet)} & 经历，体验 \\
动词 [har] & at besøge & \textit{(besøger, besøgte, har besøgt)} & 参观；拜访 \\
动词 [har] & at sidde & \textit{(sidder, sad, har siddet)} & 坐着（静态，必用 \textit{har}） \\
动词 [har] & at lægge & \textit{(lægger, lagde, har lagt)} & 放，摆（动作使平放） \\
动词 [har] & at skrive & \textit{(skriver, skrev, har skrevet)} & 写 \\
动词 [har] & at sende & \textit{(sender, sendte, har sendt)} & 发送；寄送 \\
动词 [har] & at ødelægge & \textit{(ødelægger, ødelagde, har ødelagt)} & 破坏，毁掉 \\
\midrule
名词 & et postkort & \textit{(postkortet, postkort, postkortene)} & 明信片 \\
名词 & en ferie & \textit{(ferien, ferier, ferierne)} & 假期 \\
名词 & en kirke & \textit{(kirken, kirker, kirkerne)} & 教堂 \\
名词 & en kuffert & \textit{(kufferten, kufferter, kufferterne)} & 行李箱 \\
形容词 & spændende / fantastisk & \textit{spændende (不变)} / \textit{fantastisk (-, -e)} & 令人兴奋的 / 极好的 \\
形容词 & eneste / tom & \textit{eneste (不变)} / \textit{tom (tomt, tomme)} & 唯一的（\textit{et eneste ord}）/ 空的 \\
\bottomrule
\end{tabularx}
}

% ---------------------------------------------------------------------
\ncdsubsec{四、Grammatik \& Skriveteknik（核心语法与 Modul 2 写作技巧精讲）}

\begin{grammatikbox}{语法点 1：现在完成时（Førnutid）助动词用 \textit{har} 还是 \textit{er}？}
丹麦语现在完成时的结构是：\textbf{\textit{har / er} + 动词过去分词（kort tillægsform）}。90\% 以上的动词用 \da{har}，但有两类高频动词必须用 \dabold{er}，这是 Modul 2 考试最高频的纠错点：

\vspace{0.3em}
\renewcommand{\arraystretch}{1.25}
{\small
\begin{tabularx}{\linewidth}{@{} l >{\raggedright\arraybackslash}p{4.3cm} >{\raggedright\arraybackslash}p{4.6cm} >{\raggedright\arraybackslash}X @{}}
\toprule
\textbf{助动词} & \textbf{适用条件} & \textbf{高频动词清单（务必记牢）} & \textbf{课文例句} \\
\midrule
\da{har} & \ding{172} \textbf{所有及物动词}（直接带宾语）\newline \ding{173} \textbf{静态存在/持续动作}（无 A$\rightarrow$B 位移） & \textit{være, bo, sidde, stå, ligge, arbejde, spise, drikke, se, læse, skrive, besøge, opleve, prøve, sende} & \textit{De \textbf{har været} i Italien i fem dage.}\newline \textit{Han \textbf{har siddet} på en café.} \\
\midrule
\dabold{er} & \ding{172} \textbf{不及物位移动词}（从 A 点到 B 点）\newline \ding{173} \textbf{状态突然改变}的动词 & \textit{rejse, tage afsted, ankomme, komme, gå, køre, cykle, flyve, flytte, løbe; blive, vågne, falde i søvn, dø} & \textit{Lukas og Maja \textbf{er rejst} til Rom.}\newline \textit{Det \textbf{er blevet} varmere.} \\
\bottomrule
\end{tabularx}
}
\vspace{0.3em}

\textbf{进阶微辨析（Modul 2 加分项）}：同一个位移动词，\textbf{只强调“纯动作/运动量”时用 \textit{har}，强调“到达某地/位移终点”时用 \textit{er}}！
\begin{itemize}
  \item \textit{Vi \da{har gået} meget i dag.}（今天我们走了很多路——强调动作本身，无终点 $\rightarrow$ \textit{har}）
  \item \textit{Vi \dabold{er gået} hen til museet.}（我们已经走到博物馆了——有明确终点 \textit{hen til} $\rightarrow$ \textit{er}）
\end{itemize}
\end{grammatikbox}

\begin{grammatikbox}{语法点 2：现在完成时 vs 一般过去时的“黄金切换法则”}
很多同学写作文时时态乱跳，其实丹麦语有一条极清晰的边界线：
\begin{enumerate}
  \item \textbf{用现在完成时（\textit{har/er + 分词}）的三种情况}：
    \begin{itemize}
      \item \textbf{从过去持续到现在}（常带 \textit{i fem dage, siden i mandags, indtil nu}）：\textit{Vi \da{har} allerede \da{været} her i fem dage.}（现在还在罗马！）
      \item \textbf{刚发生且影响现在}（常带 \textit{allerede, endnu, lige} 或不带具体过去时间）：\textit{Maja \da{har} allerede \da{sendt} mange billeder.}
      \item \textbf{泛泛总结到目前为止的经历}（不指明哪一天）：\textit{Vi \da{har set} mange spændende museer.}
    \end{itemize}
  \item \textbf{必须切换回一般过去时（\textit{Datid}）的情况}：只要句子里出现了\textbf{“已经彻底结束的具体过去时间点/时间段”}（\textit{i går, i mandags, den første dag, sidste år, kl. 10, før de tog afsted}），就\textbf{绝对不能用完成时}，必须立刻切成一般过去时！
\end{enumerate}
对比看 Maja 邮件中的丝滑切换：\textit{Vejret \da{har været} fantastisk, selvom det \dabold{blæste} lidt \textbf{den første dag}.}（整体天气到今天一直很好 $\rightarrow$ 完成时；“第一天”已结束 $\rightarrow$ 过去时！）
\end{grammatikbox}

\begin{grammatikbox}{写作技法：DU3 Modul 2 邮件写作（Skrivning Opgave 1）四步满分模板}
对照 Tekst B 中 Maja 的满分回信，在考场上严格执行这 4 步，即可稳拿 85--100 词（远超 70 词底线）：
\begin{enumerate}
  \item \textbf{第 1 步｜礼貌寒暄 + 接住话题（约 15 词）}：\textit{Hej [Navn]! Tak for din søde mail! Ja, det er rigtigt, at...}（用 \textit{at} 从句顺手展示语序！）
  \item \textbf{第 2 步｜逐一回答全部问题（约 55--65 词）}：检查来信有几个问号，一个都不能漏！用\textbf{完成时总结近况} + \textbf{过去时补充具体细节}，穿插 1--2 个从句连词（\textit{fordi / selvom / så}）和 1 次\textbf{时间状语提前的 V2 倒装}（如 \textit{I går gik vi...}）。
  \item \textbf{第 3 步｜互动反问 + 表达期待（约 15 词）}：\textit{Hvordan går det med dig? Jeg glæder mig til at se dig snart!}
  \item \textbf{第 4 步｜地道落款}：朋友间常用 \textit{Kærlig hilsen / Mange hilsner + 你的名字}。
\end{enumerate}
\end{grammatikbox}

% ---------------------------------------------------------------------
\ncdsubsec{五、Mønsterøvelser（关键句型操练）}

\textbf{练习 A：用 \textit{har} 或 \textit{er} + 括号内动词的正确过去分词填空}
\begin{enumerate}
  \item Lukas og Maja \underline{\hspace{1.4cm}} allerede \underline{\hspace{2.2cm}} til Italien. (\textit{at rejse})
  \item De \underline{\hspace{1.4cm}} \underline{\hspace{2.2cm}} på mange museer i Rom. (\textit{at være})
  \item Lukas \underline{\hspace{1.4cm}} endnu ikke \underline{\hspace{2.2cm}} et eneste postkort. (\textit{at skrive})
  \item Hvornår \underline{\hspace{1.4cm}} I \underline{\hspace{2.2cm}} hjem fra ferie? (\textit{at komme})
  \item I eftermiddags \underline{\hspace{1.4cm}} vi \underline{\hspace{2.2cm}} kaffe i solen. (\textit{at drikke})
  \item Min søster \underline{\hspace{1.4cm}} lige \underline{\hspace{2.2cm}} til en ny lejlighed i Aarhus. (\textit{at flytte})
  \item Jeg \underline{\hspace{1.4cm}} allerede \underline{\hspace{2.2cm}} mange billeder til min familie. (\textit{at sende})
  \item Det \underline{\hspace{1.4cm}} \underline{\hspace{2.2cm}} meget varmere i denne uge. (\textit{at blive})
\end{enumerate}

\textbf{练习 B：现在完成时 vs 一般过去时实战选填}
\begin{enumerate}
  \item Jeg \underline{\hspace{2.2cm}} (\textit{har boet / boede}) i Danmark i to år nu, og jeg er rigtig glad for det.
  \item Sidste sommer \underline{\hspace{2.2cm}} (\textit{har vi rejst / rejste vi}) til Frankrig med tog.
  \item I går \underline{\hspace{2.2cm}} (\textit{har det regnet / regnede det}) hele dagen, men i dag skinner solen.
  \item \underline{\hspace{1.4cm}} du allerede \underline{\hspace{1.8cm}} (\textit{Har ... læst / Læste}) Daniels nye e-mail?
  \item Da vi \underline{\hspace{2.2cm}} (\textit{er ankommet / ankom}) til hotellet i mandags, var vi meget trætte.
  \item Lukas \underline{\hspace{2.2cm}} (\textit{har glemt / glemte}) at købe frimærker indtil nu.
\end{enumerate}

\textbf{练习 C：句型改写（练熟 V2 倒装与从句语序）}
\begin{enumerate}
  \item 把时间状语提前到句首（注意 V2 倒装）：\textit{Vi ankom til Rom i mandags.} $\rightarrow$ \textit{I mandags...}
  \item 把句子改为否定完成时（注意 \textit{ikke} 的位置）：\textit{Han har skrevet postkortet.} $\rightarrow$
  \item 用 \textit{selvom} 连接两句话（注意从句副词位置）：\textit{Vejret har været rigtig godt. Det blæste lidt i går.} $\rightarrow$
\end{enumerate}

% ---------------------------------------------------------------------
\ncdsubsec{六、Modul 2 真题写作实战与参考答案}

\begin{emailbox}{Skrivning Opgave（DU3 Modul 2 仿真模拟题）}
\textbf{Situation}: Du har fået en mail fra din ven, Sofie.\\
\textbf{Opgave}: Læs e-mailen, og skriv et svar til Sofie. Du skal svare på alle Sofies spørgsmål. Du skal skrive \textbf{minimum 70 ord}.

\vspace{0.3em}
\textit{Hej!\\
Jeg har hørt, at du har været på weekendtur i Danmark. Det lyder hyggeligt! Hvor har du været henne? Hvem var du afsted sammen med? Hvad lavede I? Og vil du anbefale stedet til andre?\\
Mange hilsner\\
Sofie}
\end{emailbox}

\textit{（建议先独立用本课的四步模板手写回复，写完再核对下面的参考范文和练习答案！）}

\begin{answerbox}
\textbf{Før du læser 参考回答}
\begin{enumerate}
  \item \textit{De har været i Rom i fem dage, og Lukas har ikke skrevet et eneste postkort indtil nu.}
  \item \textit{I går gik de en lang tur gennem byen og spiste pizza på en lille restaurant. Vejret har været fantastisk med masser af sol, selvom det blæste lidt den første dag.}
\end{enumerate}

\textbf{练习 A 答案（\textit{har} vs \textit{er}）}
\begin{enumerate}
  \item \da{er ... rejst}（\textit{rejse} 表示位移 $\rightarrow$ \textit{er}）
  \item \da{har været}（\textit{være} 静态存在 $\rightarrow$ \textit{har}）
  \item \da{har ... skrevet}（\textit{ikke} 放在助动词与分词之间）
  \item \da{er ... kommet}（\textit{komme} 位移 $\rightarrow$ \textit{er}）
  \item \da{har ... drukket}（及物动作 $\rightarrow$ \textit{har}；注意 \textit{drikke $\rightarrow$ har drukket}）
  \item \da{er ... flyttet}（搬家属于位置/状态改变 $\rightarrow$ \textit{er}）
  \item \da{har ... sendt}（及物动作 $\rightarrow$ \textit{har}）
  \item \da{er blevet}（状态变化动词 \textit{blive} $\rightarrow$ \textit{er}）
\end{enumerate}

\textbf{练习 B 答案（完成时 vs 过去时）}
\begin{enumerate}
  \item \da{har boet}（\textit{i to år nu} 持续到现在）
  \item \da{rejste vi}（\textit{Sidste sommer} 已结束，过去时 + 句首 V2 倒装）
  \item \da{regnede det}（\textit{I går} 已结束，过去时 + 倒装）
  \item \da{Har ... læst}（\textit{allerede} 询问到现在为止的结果）
  \item \da{ankom}（\textit{i mandags} 具体过去时间点）
  \item \da{har glemt}（\textit{indtil nu} 到目前为止）
\end{enumerate}

\textbf{练习 C 答案}
\begin{enumerate}
  \item \textit{I mandags \dabold{ankom vi} til Rom.}（主句 V2 倒装：\textit{ankom} 必须在第 2 位）
  \item \textit{Han har \dabold{ikke} skrevet postkortet.}
  \item \textit{Vejret har været rigtig godt, selvom det blæste lidt i går.}（若 \textit{selvom} 从句放句首，则主句倒装：\textit{Selvom det blæste lidt i går, \dabold{har vejret} været rigtig godt.}）
\end{enumerate}

\textbf{Modul 2 仿真写作高分范文（93 词）}\\
\textit{Til: Sofie \quad|\quad Emne: Weekendtur til Skagen}

\textit{Hej Sofie}

\textit{Tak for din mail! Ja, jeg \da{har været} på weekendtur i Skagen i Nordjylland, og det \da{var} virkelig dejligt. Jeg \da{var} afsted sammen med min kæreste og to gode venner fra sprogskolen.}

\textit{Vi kørte derop fredag eftermiddag. Om lørdagen \da{gik} vi en lang tur på stranden og \da{så} Grenen, hvor de to have mødes. Bagefter spiste vi frisk fisk på havnen. Jeg vil helt klart anbefale Skagen, fordi naturen er så smuk og rolig.}

\textit{Har du selv planer om at rejse snart?}

\textit{Mange hilsner\\
\text{[Dit navn]}}

\small\textbf{点评}：精准覆盖全部 4 个问题；同时运用了 \textit{har været} 完成时、\textit{var/gik/så/spiste} 过去时、时间状语提前倒装（\textit{Om lørdagen gik vi...}）、原因从句（\textit{fordi...}）以及结尾自然反问，字数 93 词，达到 Modul 2 优秀标准。
\end{answerbox}
